\section{Related Work}
\label{sec:related}

\paragraph{Machine learning for card fraud}
Card fraud detection is a long-studied supervised learning problem with rare positives, shifting behaviour and delayed labels; recent reviews cover the breadth of methods~\cite{lucas2020credit,cherif2023credit,hafez2025systematic}. A recurring finding is that features aggregating a card's recent history improve detection substantially~\cite{whitrow2009transaction,bahnsen2016feature,jurgovsky2018sequence}. Work by Dal Pozzolo and colleagues, done with an industrial partner, set out what a realistic evaluation looks like: fraud patterns drift, labels are immediate only for the few transactions that investigators review and arrive late for the rest, and what matters operationally is the precision of the small set of alerts raised each day~\cite{dalpozzolo2014learned,dalpozzolo2018credit}. The practical handbook by Le Borgne et al.~\cite{leborgne2022fraud} turns the same principles into a recipe: training and test sets ordered in time, separated by a delay period, with prequential validation. Our chronological, gap and rolling protocols follow this line of work; our contribution is to measure what departing from it costs on two widely used public datasets.

\paragraph{Evaluation practice on public fraud datasets}
Many studies on the ULB and IEEE-CIS datasets do not evaluate this way. Random, shuffled or stratified train/test splits and shuffled $k$-fold cross-validation are common~\cite{taha2020intelligent,esenogho2022neural,alfaiz2022enhanced,chung2023credit,mim2024soft,du2024novel,becerrasuarez2026conservative}, and some pipelines resample the whole dataset before dividing it into training and test sets~\cite{ileberi2021performance,ahmad2023class,khalid2024enhancing}. The reported figures are often close to perfect, for example an accuracy of 99.98\% on ULB~\cite{ileberi2021performance}. We do not know how prevalent these practices are: we found no audit that counts them, only qualitative critiques. Hayat and Magnier~\cite{hayat2025data} review twenty methods evaluated on ULB, find pre-split preprocessing and oversampling in many of them, and show that a trivial network inside a leaky pipeline matches published results. Kabane~\cite{kabane2024impact} compares sampling before and after the split for XGBoost on ULB. Both name the missing temporal split as a problem without testing it.

\paragraph{Leakage and time-aware evaluation}
Leakage, the use of information during training that would not be available at prediction time, was formalised by Kaufman et al.~\cite{kaufman2012leakage}. Kapoor and Narayanan~\cite{kapoor2023leakage} document it across many scientific fields; their taxonomy includes both issues studied here, preprocessing or sampling before the split and test data that do not respect time. Santos et al.~\cite{santos2018cross} show that oversampling before cross-validation yields over-optimistic estimates on imbalanced data in general. In security, TESSERACT~\cite{pendlebury2019tesseract} demonstrated that malware classifiers evaluated without temporal constraints look far better than they perform on future samples, and temporal snooping is among the pitfalls catalogued by Arp et al.~\cite{arp2022dos}. Blocked and forward-chaining validation for temporally dependent data is well established in forecasting~\cite{bergmeir2012use}, and concept drift has its own literature~\cite{gama2014survey}.

\paragraph{Closest prior work}
The observation that random splits flatter fraud models is not new, and we do not claim it. The Fraud Dataset Benchmark~\cite{grover2022fraud} provides out-of-time test sets for several public datasets, including both of ours, and reports that the validation \rocauc{} of AutoML tools overstates their out-of-time \rocauc{} markedly on IEEE-CIS and little on ULB. Georgiev et al.~\cite{georgiev2026modeling} rerun one convolutional network on ULB under a chronological split in an appendix and see recall fall from 0.88 to 0.80. Breskuvien\.{e} and Dzemyda~\cite{breskuviene2024enhancing} adopt a time-based split for a feature-selection study and argue against random splitting. Zuberi et al.~\cite{zuberi2026robust} evaluate chronologically on ULB without a random arm. Hsin et al.~\cite{hsin2022feature} compare chronological and random partitioning on proprietary fund-transfer data and find the random results unrealistically high.

What is missing is a controlled, like-for-like measurement: the same models, tuning budget and training-set size under both protocols, on both public datasets, with the leak from oversampling quantified alongside, controls that separate leakage from a harder test period, and metrics tied to an alert budget. That is what this paper provides.
